\paragraph{IV 21, \S3: linear systems, fixed points and maps to projective bundles (cote25, sites 112--122).}
The handwritten heading at site 112 introduces linear systems of divisors and morphisms into projective bundles; its large section digit retains overlapping3/2 traces. The opening preliminary account of relative divisors and their representing functor is visibly canceled. The following live text defines fixed points of a family and studies the open complement of the fixed-point set. Sites113--114 compare the family on $X/S$ parametrized by $T$ with the transposed family on $T/S$, including the hypotheses on flatness and geometrically irreducible fibers, and then discuss fixed points ``au sens large''. The author separates an added sorite about the transposed family, marks its relocation to paragraph 19, and adds a qualification about its behaviour under base change. At site 115 the projective-bundle representation of the divisor functor supplies the degree of a family and the definition of a linear system. The associated morphism on the complement of the fixed points is related to the canonical incidence divisor $\underline H$ on $P\times_S\check P$. Proposition 3.1 at site 116 states the correspondence for systems without fixed points and expresses its condition through the images of associated prime cycles after extension to an algebraic closure. The handwritten additions to the preceding discussion include the non-degeneration condition and the terminology of a pseudo-morphism; the text refers back to paragraph 11 and proposes moving part of the discussion to paragraph 19.

Site 117 defines fixed components and gives a uniqueness proposition whose original typed number 2 has an unreadable handwritten replacement. Its proof compares two systems with the same pseudo-morphism, uses schematic density and replaces the earlier SGA reference with\S20.10; a marginal reference points to\S20.11. Lemma 3.2 at site 118 gives the direct-image extension statement under the printed flatness and depth conditions. Its proof reduces to the noetherian case and cites paragraphs 6 and5, with an authorial suggestion that the lemma should appear in paragraph 11. The page then completes the preceding uniqueness argument using the maximal domain of the pseudo-morphism. The long bracket and wavy pen line spanning the opening argument and lemma have unresolved editorial function, so their body text is retained. Two conspicuous printed features also remain: the quantifier in the lemma reads $x\in X$, and the displayed immersion has target $T$ in $i:X-T\to T$.

The revised proposition at site 119 returns to this correspondence using two sets, written with distinct oval-stemmed and forked-stemmed handwritten glyphs. Their alphabetic identification remains uncertain; the transcription preserves that distinction. Its heading has adjacent3-like figures and an uncertain underlying2, without a securely written separator. The statement relates systems satisfying the displayed depth condition to pseudo-morphisms and gives an invertibility criterion involving $i_*(\underline L_U)$, with $\underline L_U$ obtained from the pullback of $\underline O_P(1)$. The inserted non-degeneration condition refers explicitly to3.1. Sites120--121 continue the proof by flat base change to the dual projective bundle, tensoring with an invertible Module and comparison with the incidence divisor. The argument cites II7 and the extension lemma. Site 120's main-body formula for $D'$ lacks an inverse-image exponent, while its marginal note writes that exponent explicitly; both written forms are preserved. The long marginal note remains partly cut and partly illegible. At site 121 the first inserted lemma reference reads4.2, whereas the later references on the same page read3.2. Corollary 3.4 states the restriction and extension assertions for locally free Modules and divisors under the stated depth hypothesis, including uniqueness and schematic closure in the added note. Its clause(b) visibly says that $\underline L'$ is on $X$. Corollary 3.5 and its continuation at site 122 use local factoriality, with regularity as the example, to obtain the asserted bijection; the final handwritten qualification excludes factorization through a hyperplane after passage to $\bar k$. The proof invokes Auslander--Buchsbaum and replaces its planned ChapterIII reference with\S20. The older final paragraph at site 122 discusses an invertible Module associated with a linear system. It lies beneath crossing diagonals with an arrow-like terminal whose function is unresolved, so its words and formulas remain available. Its blank relation positions, typed superscript$x$, and uncertain decoration on the last$P$ are retained without completing the formulas.
